\subsection*{Order-sensitive reading of pictured questions}
Each 32-by-32 binary glyph was averaged into an eight-by-eight image. Two one-dimensional convolutions of width five and 32 channels processed the sequence of these images, giving a nine-glyph receptive field. For a candidate writing action, the encoder combined the mean question representation with the local representation at the first visual match for its associated field, or for its headword when the field was absent. A learned residual modified the candidate key in the existing recurrent writer. The control had identical parameters, initialization, original candidate features, targets and optimization. Its additional convolutional reader received the full glyphs sorted by their bitmap bytes. At deployment the ordered model consumes raster images, and copies historic pixels directly from the pictured card.

The two arms started from the previously trained 28,032-parameter entry writer. Eight passes over 44,600 requests gave 2,792 AdamW updates per arm. Each of the two field orders and two historical-stage orders appeared on each training headword. Three sets of status expressions supplied training supervision. Development retained the same 284 families held out from the preceding fit. Each complete response had to reproduce the reference pixels and stop autonomously. The subsequent cue-relocation test also used all 4,544 requests. A separate CPU implementation reproduced every GPU stream and stop from the ordered writer across the three registered conditions (10,224 answers). The additional ordered fit took 13.87 seconds on one RTX 4090 D with 128.43 MiB peak allocated by PyTorch. The sorted fit took 15.00 seconds. Corpus rendering, preceding training and shared preparation fall outside those fit times.
